\documentclass[11pt, a4paper, oneside]{article}
\usepackage[ngerman]{babel}
\usepackage{worksheet}

\begin{document}
	\author{L. Bung}
	\title{Exkurs: Potenzfunktionen \hspace{10cm} mit negativer Hochzahl}
	\subject{Mathematik}
	\class{FHR1}
	\maketitle
	
	Wir haben bereits Potenzfunktionen der Form $f(x) = a \cdot x^n$ mit $n \in \mathbb{N}$ kennen gelernt.
	Was passiert, wenn sogar $n \in \mathbb{Z}$ ist -- und wir also negative Zahlen einsetzen?
	
	\partnertask{Verlauf der Graphen von Potenzfunktionen mit negativer Hochzahl}
	
	a) Füllen Sie die folgende Tabelle aus:
	
	\begin{table}[H]
		\centering
		\begin{tabularx}{\textwidth}{| X | p{1.5cm} | p{1.5cm} | p{1.5cm} | p{1.5cm} | p{1.5cm} | p{1.5cm} |}
			\hline
			$x$ & -2 & -1 & 0 & 1 & 2 & 3\\
			\hline
			\hline
			$f(x) = x^{-1}$ &&&&&&\\
			\hline
			$g(x) = x^{-2}$ &&&&&&\\
			\hline
			$h(x) = x^{-3}$ &&&&&&\\
			\hline
			$i(x) = x^{-4}$ &&&&&&\\
			\hline
		\end{tabularx}
	\end{table}
	
	b) Zeichnen Sie die Graphen der Funktionen $f(x)$, $g(x)$, $h(x)$ und $i(x)$ in ein Koordinatensystem.
	Verwenden Sie dazu die Tabelle aus Teilaufgabe a).
	
	\checkered[10cm]
	
	c) Beschreiben Sie, wie die Graphen der Funktionen verlaufen.
	Welche Muster fallen Ihnen auf?
	Wie verlaufen wohl die Graphen der Funktionen $j(x) = x^{-5}$ und $k(x) = x^{-6}$?
	
	\lines
	
	\begin{hint}{Verlauf der Graphen von Potenzfunktionen mit negativer Hochzahl}
		\phantom{x}\vspace{6.5cm}
	\end{hint}
	
	\singletask{Funktionen mit Asymptoten bestimmen}
	
	Geben Sie jeweils eine Funktion an, die die folgenden Asymptoten hat:
	
	\begin{multicols}{3}
		\begin{enumerate}[label=\alph*)]
			\item $x=0$ und $y=2$
			\item $x=-3$ und $y=0$
			\item $x=1$ und $y=-1$
		\end{enumerate}
	\end{multicols}
	
	\checkered[4cm]
	
	\let\clearpage\relax
\end{document}